\documentclass{article}

\title{Ch. 31 CW \#2 Solutions}
\author{Principles of Macroeconomics}
\date{\today}

\begin{document}

\maketitle

\begin{enumerate}

\item 

	\begin{enumerate}
	
	\item People will need more money for their holiday shopping, so their preference for liquidity will increase. That will shift demand for money to the right.
	
	\item Since people can now buy fast food on credit, they will need less money than they used to. Their preference for liquidity will decrease, so their demand for money will shift to the left.
	
	\item An open-market purchase will shift the money supply to the right. That will cause a movement along the demand curve to the right.
	
	\end{enumerate}
	
\item Recall that the quantity equation is \[MV = PY\] where $PY$ is nominal GDP. In all four years, nominal GDP is growing slower than the money supply. For the quantity equation to hold, the velocity of money must have been growing at a negative rate.

\item

	\begin{enumerate}
	
	\item Recall that $MV = PY$. If $M$ and $V$ are constant and $Y$ is shrinking, then $P$ must increase to offset the change in $Y$. 
	
	\item If $V$ and $P$ are constant, then $M$ must shrink at the same rate as $Y$.
	
	\end{enumerate}
	
\item

	\begin{enumerate}
	
	\item It helps Lynn, and it hurts Boris. If inflation is actually 4\%, then the real interest rate on the loan is actually $8\% - 4\% = 4\%$. That means Lynn receives more and Boris pays more than expected in terms of purchasing power.
	
	\item It hurts Lynn, and it helps Boris. If inflation is actually 7\%, then the real interest rate on the loan is actually $8\% - 7\% = 1\%$. That means Lynn receives less and Boris pays less than expected in terms of purchasing power.
	
	\end{enumerate}

\item

	\begin{enumerate}
	
	\item There is a net cost on the economy because it disincentivizes people from holding money. That creates shoe leather costs.
	
	\item There is no net cost on the economy, but it does arbitrarily redistribute wealth from the lender to Hector.
	
	\item There is a net cost on the economy from setting and advertising new prices. That creates menu costs. 
	
	\end{enumerate}
	
\item They would have been especially attractive between 1993 and 1998 because the nominal interest rate was below the real inflation rate. That means real interest rates were negative, so borrowing was profitable to borrowers, not lenders.

\end{enumerate}

\end{document}